%! TeX root = ../main.tex

\section{Conclusiones}

En este proyecto se realizó el análisis de un motor de DC a fin de obtener
información sobre su estabilidad, controlabilidad y observabilidad. El modelado
fue la parte más interesante, puesto que la consideración de la velocidad angular
jugaba un papel bastante importante en la forma en que el sistema en espacio 
de estados tendería a comportase en el análisis, cambiando sus características
de estabilidad.

Otro punto importante a resaltar es la complejidad que puede surgir de un sistema
tan simple como un motor de DC al momento de realizar los análisis, lo que permite
vislumbrar un panorama importante al momento de considerar sistemas con procesos
más complejos como un avión, procesos poblacionales de animales, o incluso este
mismo motor, pero con carga.

El análisis mostró que este sistema es marginalmente estable debido a la posición
angular, la cual no afecta a nuestra salida de interés (la velocidad angular) y que
se mantiene creciente en el tiempo sin llegar a un punto estable. Cuestión corroborada
con los autovalores de la matriz dinámica y el análisis de Lyapunov.

El análisis numérico indica que el sistema es controlable. Sin embargo, debido a esta variable de posición angular,
el sistema es no observable.
